% Part: sets-functions-relations
% Chapter: functions
% Section: partial-functions

\documentclass[../../../include/open-logic-section]{subfiles}


\begin{document}

\olfileid{sfr}{fun}{par}

\olsection{Partiële functies}

\begin{explain}
Soms is het nuttig de definitie van een functie te versoepelen, zodat
de functiewaarde niet voor elke mogelijke invoer gedefinieerd hoeft te
zijn. Zulke afbeeldingen heten \emph{partiële functies}.
\end{explain}

\begin{defn}
Een \emph{partiële functie} $f \colon A \pto B$ is een afbeelding die
aan elk !!{element} van~$A$ ten hoogste één !!{element} van~$B$
toekent. Als $f$ een element van~$B$ aan $x \in A$ toekent, noemen
we $f(x)$ \emph{gedefinieerd}; anders noemen we dit
\emph{ongedefinieerd}. Als $f(x)$ gedefinieerd is, schrijven we $f(x)
\fdefined$; anders schrijven we $f(x) \fundefined$. Het
\emph{domein} van een partiële functie~$f$ is de deelverzameling
van~$A$ waarop ze gedefinieerd is, dus $\dom{f} = \Setabs{x \in
A}{f(x) \fdefined}$.
\end{defn}

\begin{ex}
Elke functie $f\colon A \to B$ is ook een partiële functie. Partiële
functies die overal op~$A$ gedefinieerd zijn---dus wat we tot nu toe
kortweg een functie noemden---heten ook \emph{totale} functies.
\end{ex}

\begin{ex}
De partiële functie $f \colon \Real \pto \Real$ gegeven door $f(x) =
1/x$ is ongedefinieerd voor $x = 0$ en overal elders gedefinieerd.
\end{ex}

\begin{prob}
Gegeven $f\colon A \pto B$, definieer de partiële functie $g\colon B
\pto A$ als volgt. Voor elke $y \in B$: als er precies één $x \in A$
is waarvoor $f(x) = y$, stel dan $g(y) = x$; schrijf anders $g(y)
\fundefined$. Toon aan dat als $f$ injectief is, $g(f(x)) = x$ voor
alle $x \in \dom{f}$ en $f(g(y)) = y$ voor alle $y \in \ran{f}$.
\end{prob}

\begin{defn}[Grafiek van een partiële functie]
Zij $f\colon A \pto B$ een partiële functie. De \emph{grafiek} van~$f$
is de relatie $R_f \subseteq A \times B$ die wordt gedefinieerd door
\[
R_f = \Setabs{\tuple{x,y}}{f(x) = y}.
\]
\end{defn}

\begin{prop}
Stel dat $R \subseteq A \times B$ de eigenschap heeft dat uit $Rxy$ en
$Rxy'$ volgt dat $y = y'$. Dan is $R$ de grafiek van de partiële
functie $f\colon A \pto B$ die als volgt wordt gedefinieerd: als er een
$y$ is waarvoor $Rxy$, stel dan $f(x) = y$; schrijf anders $f(x)
\fundefined$. Als $R$ bovendien \emph{serieel} is, dat wil zeggen dat
er voor elke $x \in A$ een $y \in B$ bestaat waarvoor $Rxy$, dan is
$f$ totaal.
\end{prop}

\begin{proof}
Stel dat er een $y$ is waarvoor $Rxy$. Als er nog een $y' \neq y$ was
waarvoor $Rxy'$, zou de voorwaarde voor $R$ worden geschonden. Als er
een $y$ is waarvoor $Rxy$, is die $y$ dus uniek en is $f$
welgedefinieerd. Verder geldt onmiddellijk $R_f = R$, en $f$ is totaal
als~$R$ serieel is.
\end{proof}
\end{document}
